% RESEARCH_DRAFT_SHA256: f67fcb59b75a152758585a057505795d7613210e012be0c1b6f646bcaaf89ae0
% Layout and source-reference wrappers expanded for the web renderer.
\documentclass[11pt]{article}
\usepackage[margin=25mm]{geometry}
\usepackage{fontspec,unicode-math}
\setmainfont{Latin Modern Roman}
\setmathfont{Latin Modern Math}
\usepackage{amsmath,amssymb,mathtools,xurl,hyperref}
\hypersetup{colorlinks=true,urlcolor=blue,linkcolor=blue}
\setcounter{secnumdepth}{0}
\setlength{\parindent}{0pt}
\setlength{\parskip}{5pt}
\setlength{\emergencystretch}{3em}
\newcommand{\sref}[2]{\hyperlink{source:#1:#2}{[S#2]}}
\newcommand{\cataloguescope}[1]{\paragraph{Source record.}#1}
\newcommand{\floor}[1]{\left\lfloor#1\right\rfloor}
\begin{document}
\section*{Research-attempt provenance}
This standalone research attempt retains the complete original definition body below,
including its historical source and preparation notes. Those notes describe the
earlier catalogue preparation; the new attempt and its verification are identified
separately after that body. No original ID, catalogue status, or source statement is changed.

% UnsolvedMath ID 20001239; source catalogue code AIM-COMPUTATION-0077
\section{A scale-normalized robustness criterion for algebraic design fans}\label{20001239}
{\small UnsolvedMath ID 20001239 · Computer Science}
\subsection{Definitions and mathematical statement}

\paragraph{Shared notation.}Standard number systems are $\mathbb N,\mathbb Z,\mathbb Q,\mathbb R,\mathbb C$. Specialized notation and index ranges follow the cited source and the definitions below.


A statistical design specifies experimental settings, often represented by a finite point set. A design perturbation changes those settings. An estimable linear parameter functional is one determined by the model's expected observations; comparing estimators normally requires a loss or variance criterion.
\paragraph{Record scope and qualifications.}
The catalogue title is retained as metadata. Named certificates, contracts, bounds or frameworks in that title are not asserted as results of the original source extract.
The extracted term Est tau factor is undefined, as are the perturbation law and meaning of best. No fan, robustness threshold, or scale-normalized criterion is reconstructed. The supplied catalogue marks this record partially solved; its present research status was not independently checked.
The source text below is rendered literally where necessary; damaged formulas and composite extracts are not silently repaired into a new mathematical claim.
\paragraph{Statement recorded in the supplied source.}
Question 62 (Riccomagno). Is there an algebraic method to decide which Est τ factor is best when working with perturbed designs? \href{https://aimath.org/WWN/compalgstat/compalgstat.pdf}{[S2]}
\paragraph{Source status.}
The supplied catalogue labels this source record partially solved; that metadata was checked against the supplied snapshot on 2026-10-11. The label is preserved for provenance and is not an independent verification of present open status or mathematical correctness.
\subsection{Short English statement}
The source asks about algebraic choices of estimators or factors under perturbed designs.
\subsection{Sources}
\begin{itemize}
\item[S1] Ulam AI and the UnsolvedMath Contributors, UnsolvedMath: A Curated Collection of Open Mathematics Problems, record 20001239 (AIM-COMPUTATION-0077). Catalogue curation and metadata: CC BY 4.0. The source extract is retained; definitions, English summaries and qualifications are editorial additions.
\url{https://huggingface.co/datasets/ulamai/UnsolvedMath}.
\item[S2] AIM Workshop Problem Lists, source linked by catalogue record AIM-COMPUTATION-0077.
\url{https://aimath.org/WWN/compalgstat/compalgstat.pdf}.
\end{itemize}
\paragraph{Preparation and review.}
The original source and UnsolvedMath contributors retain attribution for the source extract. Codex (GPT-6; exact serving revision not exposed) assisted with definitions, English summaries, scope notes and formatting. No new human mathematical authorship or independently verified solution is asserted. Permanent source ID 20001239 and code AIM-COMPUTATION-0077 checked against the complete supplied registry and current repository. Distinct upstream records, including repeated questions and broad or damaged extracts, are retained for complete coverage of the source CS category. A repeated record is not a new independent conjecture. Complete mathematical-equivalence and current-literature review has not been performed. Preparation and final review follow contribution prompts C and D, with source gaps disclosed. Structural and rendering checks do not establish a new theorem or open status.
\label{end:20001239}

\section{Literature and context check}
The primary source defines standard monomials using a term order, then asks Question 62 without a perturbation law or a definition of best. A concrete interpolation-conditioning route is tested below. It is not an identification of the unspecified preferred factor.
\section{Research attempt: Basis dependence and a normalized perturbation certificate}
For a square invertible evaluation matrix $V$ of a chosen polynomial basis at a design, the raw Euclidean condition number depends on basis scaling. For example $V=I_2$ has condition number one, while replacing its basis by the same two functions scaled by $K$ and $1$ gives $VS=\operatorname{diag}(K,1)$, with condition number $K$ for $K\ge1$. The interpolation space and its predictions are unchanged after the inverse coefficient change. Thus minimizing a raw condition number across arbitrarily normalized bases is not an intrinsic design criterion.

A comparison that removes this particular ambiguity uses a perturbed evaluation matrix $V'=V+E$ for the same basis, and the transfer matrix
\[
T=V'V^{-1}=I+EV^{-1}.
\]
It maps the vector of values on the original design to the interpolated values on the perturbed design. An invertible basis change replaces $(V',V)$ by $(V'S,VS)$ and leaves $T$ unchanged, since $(VS)^{-1}=S^{-1}V^{-1}$. If $\|EV^{-1}\|_2=\eta<1$, the triangle and reverse triangle inequalities give
\[
(1-\eta)\|u\|_2\le\|Tu\|_2\le(1+\eta)\|u\|_2.
\]
The lower bound gives injectivity; since $T$ is square, it is invertible. Applying the lower bound to $T^{-1}v$ gives $\|T^{-1}\|_2\le(1-\eta)^{-1}$, and hence $\kappa_2(T)\le(1+\eta)/(1-\eta)$. This is a proved certificate conditional on an actual perturbation bound, rather than an order-selection rule.

Small coordinate perturbations alone cannot supply a uniform certificate. With the basis $\{1,x\}$ and design $\{0,\varepsilon\}$,
\[
V=\begin{pmatrix}1&0\\1&\varepsilon\end{pmatrix},\quad
V'=\begin{pmatrix}1&0\\1&\varepsilon+\delta\end{pmatrix},\quad
T=\begin{pmatrix}1&0\\-\delta/\varepsilon&1+\delta/\varepsilon\end{pmatrix}.
\]
At $\delta=-\varepsilon$ the perturbed points coincide and $T$ is singular. The physical change $|\delta|$ can be arbitrarily small. Even if coincident designs are excluded, $\delta$ arbitrarily close to $-\varepsilon$ makes the inverse unbounded: the determinant $1+\delta/\varepsilon$ tends to zero while $\|T\|_2\ge1$. Thus separation or an appropriate relative perturbation restriction is indispensable for this route.
\section{Verification and adversarial audit}
Matrix multiplication checks the displayed transfer matrix. The singular boundary has $\eta\ge1$, so it does not contradict the strict certificate hypothesis. Basis invariance is only invariance of this transfer matrix under simultaneous basis change; it does not assert invariance under a change of design metric.
\paragraph{Audit corrections.} Raw coefficient conditioning is separated from prediction transfer. A coordinate-small perturbation is not silently assumed small after multiplication by $V^{-1}$.
\paragraph{External computation.} None; symbolic arguments and boundary checks only.
\section{Scope and remaining gaps}
The source's preferred factor, loss, and perturbation model remain undefined. The result is a standard interpolation stability calculation and an exact obstruction to an unrestricted criterion, with no new original-target theorem.
\paragraph{Integration review.} Codex (GPT-6; exact serving revision not exposed) reviewed the written arguments and expanded layout and source-reference wrappers. Repository metadata and the production web renderer were checked using Python 3.13 and JavaScript with MathJax 3.2.2. This model-assisted review is not independent human proof verification; no PDF compilation or formal proof checking is asserted.

\section{Completion Estimate}
\noindent\textbf{COMPLETION: 0\%}
The calculations below the preserved source statement investigate a concrete route,
but establish no substantiated new progress toward the entire original target.
Known special cases and failed sufficient conditions are not credited as a solution.
\section{References}
The original source references and catalogue attribution are preserved in the
Sources section above. The following URLs were actually consulted for this attempt;
retention of an additional original reference does not claim fresh consultation of it.
\begin{itemize}
\item \url{https://aimath.org/WWN/compalgstat/compalgstat.pdf} (consulted 2026-10-11). 
\end{itemize}
\section{Model and contributor disclosure}
OpenAI gpt-6-astra, requested reasoning configuration ultra, generated this research
attempt on 2026-10-11. The exact serving revision is not exposed. Codex GPT-6
performed orchestration, repository-format review, and a model-assisted mathematical
audit of the written arguments; its exact serving revision
is also not exposed. The mathematical route, boundary checks and final adversarial
self-audit followed contribution prompts A then B in
\texttt{docs/contribution\_prompts.md}, after reading \texttt{CONTRIBUTING.md}.
Python 3.14.2 assembled the standalone documents; any mathematical computation is
specified in Verification. No theorem prover was used. No human mathematical
authorship, substantial human proof correction, or independent human proof checking
is asserted. Model-assisted review, source retrieval, and structural checks do not
establish correctness or novelty. Full private reasoning transcripts are unavailable
as reproducibility artifacts; the complete arguments, inputs and outcomes actually
used are stated here or in the companion batch audit.
\section{Final status}
Unresolved as a new research contribution. This is a concrete route investigation,
not a claim of new complete resolution or a current-status certification. A known
answer reproduced for context does not become a new contribution. The catalogue's
source-status label is preserved; this attempt status does not assert that a
source-labelled solved question is open.
\end{document}
